\subsection{Experimental Design}
\label{subsec:exp-design}

\subsubsection{Test Targets}
\label{subsubsec:targets}

We evaluate GMPFuzz on three widely-deployed open-source MQTT brokers that
differ in implementation language, architecture, and target deployment
scenario.  Selecting multiple brokers avoids evaluation bias toward a single
codebase and demonstrates the protocol-level generality of our approach.

\begin{table}[htbp]
  \centering
  \caption{MQTT Broker Test Targets}
  \label{tab:targets}
  \renewcommand{\arraystretch}{1.25}
  \begin{tabular}{llllp{4.2cm}}
    \toprule
    \textbf{Target} & \textbf{Version} & \textbf{Language} & \textbf{License} &
    \textbf{Characteristics} \\
    \midrule
    Mosquitto~\cite{mosquitto}
      & v1.5.5
      & C
      & EPL-2.0
      & Lightweight reference broker; widely used in IoT devices and embedded
        systems; extensive state-machine for QoS 1/2 and retained messages. \\
    \addlinespace
    Mongoose~\cite{mongoose}
      & v7.20
      & C
      & GPLv2 / commercial
      & Event-driven single-thread network library with an embedded MQTT broker;
        compact code base suited for resource-constrained environments. \\
    \addlinespace
    NanoMQ~\cite{nanomq}
      & v0.21.10
      & C
      & MIT
      & High-performance multi-threaded broker targeting edge computing;
        implements MQTT 3.1.1 with NNG-based asynchronous I/O. \\
    \bottomrule
  \end{tabular}
\end{table}

All three brokers implement MQTT 3.1.1~\cite{mqtt311spec} and are compiled
with AFL-Clang-Fast instrumentation (edge-coverage feedback) and AddressSanitizer
(\texttt{AFL\_USE\_ASAN=1}) to amplify crash detection.  A separate
coverage-only build (GCC with \texttt{-fprofile-arcs -ftest-coverage -O0}) is
maintained for post-run branch-coverage measurement via \texttt{gcovr}.

\subsubsection{Baselines and Benchmarks}
\label{subsubsec:baselines}

We compare GMPFuzz against three representative protocol fuzzers that serve as
natural baselines for MQTT fuzzing.
\textbf{AFLNet}~\cite{aflnet} is a greybox network fuzzer that extends AFL
with a response-guided state machine.  It is the most widely adopted
protocol fuzzer and the direct foundation of GMPFuzz's fuzzing engine;
we run four independent AFLNet instances per target to match GMPFuzz's
internal parallelism of four state pools.
\textbf{MPFuzz}~\cite{mpfuzz} is a multi-client MQTT fuzzer that
coordinates publish/subscribe client pairs to trigger cross-client protocol
bugs; it executes four internal agents concurrently, providing a
competitive parallel baseline.
\textbf{Peach}~\cite{peach} is a generation-based fuzzer driven by
hand-crafted MQTT pit files, representing the class of specification-guided
fuzzers that do not leverage runtime feedback; we run four independent
instances to match the parallelism of the other tools.

All tools are containerised via Docker and executed under identical
network and process conditions.  Fuzzer binaries for AFLNet are taken from the
ProFuzzBench~\cite{profuzzbench} framework to ensure reproducibility; MPFuzz
and Peach use the images provided in the respective benchmark sub-directories
of our artefact.

\subsubsection{Hardware and Software Environment}
\label{subsubsec:env}

All experiments are conducted on a single workstation with the following
configuration.

\paragraph{Hardware.}
The host machine is equipped with an Intel Xeon W-3400 series processor
(24 physical cores), 128\,GB of DDR5 ECC memory, and an NVMe SSD
(sequential read $>$\,3\,GB/s).  All fuzzing traffic is confined to the
loopback interface; no external network connectivity is used during
experiments.

\paragraph{Software.}
The host runs Ubuntu 20.04.6 LTS (kernel 5.15).  Each fuzzer instance
executes inside an isolated Docker 24.0 container that shares the host
network namespace to minimise TCP round-trip latency.  Instrumented target
binaries are compiled with LLVM/Clang~12 (	exttt{afl-clang-fast}),
while coverage builds use GCC~9.4 with 	exttt{-fprofile-arcs
-ftest-coverage}.  LLM inference is provided by CodeLlama-13b-hf
deployed via vLLM~\cite{vllm} on a dedicated GPU node (NVIDIA A100
40\,GB), keeping LLM latency off the fuzzing host's CPU budget.
The GMPFuzz orchestrator and all analysis scripts run under Python~3.10.

\paragraph{Experimental parameters.}
The full-mode GMPFuzz run uses a fuzz-only time budget of $T_{\text{budget}} =
86{,}400$\,s (24\,h) per target, with per-epoch bounds $T_{\min} = 1{,}800$\,s
and $T_{\max} = 7{,}200$\,s.  Four state pools run in parallel, each driving
an independent AFL-Net container.  The LLM variant-generation phase
($T_{\text{LLM}} \approx 2{,}400$\,s per generation) is excluded from the
fuzz budget and runs sequentially between epochs.  All competing tools run for
the same wall-clock budget of 86,400\,s.  Each configuration is repeated once;
we report edge-coverage curves and final edge counts as primary metrics.